% GENERATED FILE. Edit canonical lessons, then run npm run generate:books.
% canonical-source-hash: fnv1a64:eda0e839a6bb4c6b
% canonical-lessons: PT-C417-molhe, PT-C417-pontao, PT-C417-embarcar, PT-C417-desembarcar, PT-C417-zarpar

\chapter{No porto}
\label{ch:pt-no-porto}
\hlchaptermodality{417}

\begin{chapteropening}
\textbf{By the end of this chapter:} \emph{I can get around a harbour: the breakwater, the pontoon, boarding, going ashore and setting sail.}

Complete the last lesson of chapter 417: I can get around a harbour: the breakwater, the pontoon, boarding, going ashore and setting sail.
\end{chapteropening}

\section[o molhe]{o molhe — the breakwater, the pier}

\label{lesson:PT-C417-molhe}



\begin{quote}
Before the new one: say the Portuguese for to cancel, then the Portuguese for the change.
\end{quote}

\subsection*{You'll want to know: o molhe}
\textbf{o molhe} — \textquotedblleft{}the breakwater, the pier\textquotedblright{}.

\begin{quote}
Ao fim da tarde, há pescadores no molhe.
\end{quote}

\begin{quote}
\emph{In the late afternoon, there are fishermen on the breakwater.}
\end{quote}

A long stone wall that protects a harbour from the waves.

\begin{cousinweb}
\emph{molhe} came through Catalan \emph{moll} 'mole, jetty' (EtymDB) from Latin \emph{mōlēs} 'mass, great stone structure' (Lewis \& Short).
\end{cousinweb}

\begin{grammarlens}[title={What you've built}]
One more word for telling, explaining and describing at B1.
\end{grammarlens}

\subsection*{Your turn}
\begin{itemize}
\raggedright
  \item \emph{Say it:} \emph{o molhe}
  \item \emph{Say it:} \emph{o molhe}, once more
  \item \emph{Say it:} a sentence of your own with \emph{o molhe}
  \item \emph{Recall:} say the Portuguese for to cancel, then the Portuguese for the change, then say \emph{o molhe} again
\end{itemize}

\begin{tcolorbox}[breakable,colback=teal!4,colframe=teal!35!black,title={Before you move on}]
What does o molhe mean? (\textquotedblleft{}The breakwater, the pier\textquotedblright{}.) Say it once more.
\end{tcolorbox}
\section[o pontão]{o pontão — the pontoon, the jetty}

\label{lesson:PT-C417-pontao}



\begin{quote}
Before the new one: say the Portuguese for the change, then the Portuguese for the breakwater.
\end{quote}

\subsection*{You'll want to know: o pontão}
\textbf{o pontão} — \textquotedblleft{}the pontoon, the jetty\textquotedblright{}.

\begin{quote}
O barco pequeno está preso ao pontão.
\end{quote}

\begin{quote}
\emph{The small boat is tied to the pontoon.}
\end{quote}

A floating platform where small boats stop.

\begin{cousinweb}
\emph{pontão} is inherited through Old Portuguese \emph{ponton} from Latin \emph{pontō} 'ferryboat, pontoon', built on \emph{pōns, pontis} 'bridge' (EtymDB; Lewis \& Short). EtymDB also reads it as \emph{ponte} with \emph{-ão}.
\end{cousinweb}

\begin{grammarlens}[title={What you've built}]
One more word for telling, explaining and describing at B1.
\end{grammarlens}

\subsection*{Your turn}
\begin{itemize}
\raggedright
  \item \emph{Say it:} \emph{o pontão}
  \item \emph{Say it:} \emph{o pontão}, once more
  \item \emph{Say it:} a sentence of your own with \emph{o pontão}
  \item \emph{Recall:} say the Portuguese for the change, then the Portuguese for the breakwater, then say \emph{o pontão} again
\end{itemize}

\begin{tcolorbox}[breakable,colback=teal!4,colframe=teal!35!black,title={Before you move on}]
What does o pontão mean? (\textquotedblleft{}The pontoon, the jetty\textquotedblright{}.) Say it once more.
\end{tcolorbox}
\section[embarcar]{embarcar — to board, to embark}

\label{lesson:PT-C417-embarcar}



\begin{quote}
Before the new one: say the Portuguese for the breakwater, then the Portuguese for the pontoon.
\end{quote}

\subsection*{You'll want to know: embarcar}
\textbf{embarcar} — \textquotedblleft{}to board, to embark\textquotedblright{}.

\begin{quote}
Os passageiros embarcam às nove.
\end{quote}

\begin{quote}
\emph{The passengers board at nine.}
\end{quote}

Also for planes and trains: \emph{embarcar no avião}.

\begin{cousinweb}
\emph{embarcar} is built from \emph{em-} 'into' and \emph{barca} 'boat', inherited from Late Latin \emph{barca} (Lewis \& Short).
\end{cousinweb}

\begin{grammarlens}[title={What you've built}]
One more word for telling, explaining and describing at B1.
\end{grammarlens}

\subsection*{Your turn}
\begin{itemize}
\raggedright
  \item \emph{Say it:} \emph{embarcar}
  \item \emph{Say it:} \emph{embarcar}, once more
  \item \emph{Say it:} a sentence of your own with \emph{embarcar}
  \item \emph{Recall:} say the Portuguese for the breakwater, then the Portuguese for the pontoon, then say \emph{embarcar} again
\end{itemize}

\begin{tcolorbox}[breakable,colback=teal!4,colframe=teal!35!black,title={Before you move on}]
What does embarcar mean? (\textquotedblleft{}To board, to embark\textquotedblright{}.) Say it once more.
\end{tcolorbox}
\section[desembarcar]{desembarcar — to go ashore; to get off, to land}

\label{lesson:PT-C417-desembarcar}



\begin{quote}
Before the new one: say the Portuguese for the pontoon, then the Portuguese for to board.
\end{quote}

\subsection*{You'll want to know: desembarcar}
\textbf{desembarcar} — \textquotedblleft{}to go ashore; to get off, to land\textquotedblright{}.

\begin{quote}
Desembarcámos numa ilha pequena.
\end{quote}

\begin{quote}
\emph{We went ashore on a small island.}
\end{quote}

Also for planes: \emph{desembarcar no aeroporto}.

\begin{cousinweb}
\emph{desembarcar} is built from \emph{des-} and \emph{embarcar} 'to board', on \emph{barca} 'boat', inherited from Late Latin \emph{barca} (Lewis \& Short).
\end{cousinweb}

\begin{grammarlens}[title={What you've built}]
One more word for telling, explaining and describing at B1.
\end{grammarlens}

\subsection*{Your turn}
\begin{itemize}
\raggedright
  \item \emph{Say it:} \emph{desembarcar}
  \item \emph{Say it:} \emph{desembarcar}, once more
  \item \emph{Say it:} a sentence of your own with \emph{desembarcar}
  \item \emph{Recall:} say the Portuguese for the pontoon, then the Portuguese for to board, then say \emph{desembarcar} again
\end{itemize}

\begin{tcolorbox}[breakable,colback=teal!4,colframe=teal!35!black,title={Before you move on}]
What does desembarcar mean? (\textquotedblleft{}To go ashore; to get off, to land\textquotedblright{}.) Say it once more.
\end{tcolorbox}
\section[zarpar]{zarpar — to set sail, to weigh anchor}

\label{lesson:PT-C417-zarpar}



\begin{quote}
Before the new one: say the Portuguese for to board, then the Portuguese for to go ashore.
\end{quote}

\subsection*{You'll want to know: zarpar}
\textbf{zarpar} — \textquotedblleft{}to set sail, to weigh anchor\textquotedblright{}.

\begin{quote}
O navio zarpou ao nascer do sol.
\end{quote}

\begin{quote}
\emph{The ship set sail at sunrise.}
\end{quote}

Informal: to leave in a hurry. \emph{Ele zarpou logo} — he was off at once.

\begin{cousinweb}
\emph{zarpar} was borrowed from Spanish \emph{zarpar}, which came from Old Italian \emph{sarpare} 'to weigh anchor' (EtymDB).
\end{cousinweb}

\begin{grammarlens}[title={What you've built}]
One more word for telling, explaining and describing at B1.
\end{grammarlens}

\subsection*{Your turn}
\begin{itemize}
\raggedright
  \item \emph{Say it:} \emph{zarpar}
  \item \emph{Say it:} \emph{zarpar}, once more
  \item \emph{Say it:} a sentence of your own with \emph{zarpar}
  \item \emph{Recall:} say the Portuguese for to board, then the Portuguese for to go ashore, then say \emph{zarpar} again
\end{itemize}

\begin{tcolorbox}[breakable,colback=teal!4,colframe=teal!35!black,title={Before you move on}]
What does zarpar mean? (\textquotedblleft{}To set sail, to weigh anchor\textquotedblright{}.) Say it once more.
\end{tcolorbox}
